% Requires: \usepackage{booktabs}
% Auto-derived from analysis/results/RQ1/rq1_filtering/rq1_filtering.md
% (regenerated by analysis/notebooks_v2/rq1_filtering.py).

\section{Pre-analysis filtering: excluded behaviour, excluded scheme, and
coherence gate}
\label{app:naive}

Three a-priori filters are applied before the RQ1 statistical analysis. They are
\emph{different kinds} of filter---the first removes a behaviour for trait
validity, the second removes an entire injection scheme for an angle-dependent
dose confound, and the third is a row-level coherence gate---and are motivated
separately below. The filters are applied in this order so that each operates on
the output of the previous one: the scheme confound is established on the
\emph{raw} (un-gated) coherence, exactly the evidence it requires, and the
row-level gate is then applied only to the surviving valid schemes. All numbers
are read from the consolidated frames at coherence column \texttt{coh\_steered};
the high positive-cosine subset is $c_{ij}>0.25$.

The behaviour \texttt{power\_seeking} appears only in the \texttt{normFalse}
scheme, which is exactly why that scheme carries $36$ raw pairs while the other
two carry $28$ (Table~\ref{tab:pair-inventory}).

\begin{table}[h]
  \centering
  \small
  \begin{tabular}{lccc}
    \toprule
    scheme & raw pairs & \texttt{power\_seeking} pairs & after filter \\
    \midrule
    \texttt{normFalse} & 36 & 8 & 28 \\
    \texttt{normTrue}  & 28 & 0 & 28 \\
    \texttt{per\_axis} & 28 & 0 & 28 \\
    \bottomrule
  \end{tabular}
  \caption{Pair inventory by injection scheme. The eight \texttt{power\_seeking}
  pairs present only in \texttt{normFalse} account for the $36$-vs-$28$
  difference.}
  \label{tab:pair-inventory}
\end{table}

\paragraph{Filter 1 (behaviour): drop \texttt{power\_seeking} for trait
validity, not coherence.}
\texttt{power\_seeking} was dropped a priori at the design commit for the
normalized re-run. The reason is that its steering vector is
\emph{non-expressing}: steering toward \texttt{power\_seeking} does not even
produce the behaviour. Over the eight pairs that contain it (all in
\texttt{normFalse}), the \texttt{power\_seeking} axis has a \emph{negative} mean
single-trait effect ($\Delta=-9.6$), and stays negative under joint steering
($\Delta=-9.0$), whereas its partner behaviour expresses normally
(Table~\ref{tab:powerseeking}). This matches the earlier single-trait
validation, in which the vector goes negative under increasing dose---a
context-conditioning failure rather than a usable direction.

\begin{table}[h]
  \centering
  \small
  \begin{tabular}{lcc}
    \toprule
    axis & mean single-trait $\Delta$ & mean joint $\Delta$ \\
    \midrule
    \texttt{power\_seeking} (dropped axis) & $-9.6$ & $-9.0$ \\
    partner behaviour                      & $45.6$ & $26.1$ \\
    \bottomrule
  \end{tabular}
  \caption{Single- and joint-steering trait effects for the eight
  \texttt{power\_seeking} pairs. The dropped axis fails to express even on its
  own; its partner is unaffected by the defect.}
  \label{tab:powerseeking}
\end{table}

Consequently every pair containing \texttt{power\_seeking} has a structurally
dead axis: the ``did axis $a$ survive joint steering?'' question is undefined,
and such a pair can never exhibit genuine emergence. Its observed regimes are
limited to mixed ($\times 5$) and additive ($\times 3$). Crucially, this is
\emph{not} a coherence filter---the \texttt{power\_seeking} pairs have high
coherence (mean $86.4$). They are removed for trait validity, not fluency.

\paragraph{Filter 2 (scheme): drop the unnormalized \texttt{normFalse} sum for
an angle-dependent dose confound.}
The \texttt{normFalse} scheme is the earlier unnormalized-sum baseline, which
injects $\alpha(\hat v_i+\hat v_j)$. Its total perturbation norm is
$\alpha\sqrt{2+2c_{ij}}$, so high-cosine pairs receive a larger total
perturbation (up to ${\sim}85\%$ more than either single condition). This
produces coherence collapse in the high-cosine, over-dosed pairs and creates a
misleading relationship between geometry and the discrete composition regimes.
Because the degradation is correlated with $c_{ij}$---the very predictor of the
RQ1 analysis---it \emph{biases} the regime comparison rather than merely adding
noise. The confound is therefore established on the \emph{raw} (un-gated)
coherence, which is the evidence it requires.

The three coherence tables below are computed on \texttt{power\_seeking}-excluded
rows, so every scheme is compared on the same $28$-pair design (apples to
apples). Removing the high-coherence \texttt{power\_seeking} pairs also
un-inflates the \texttt{normFalse} mean.

On the raw (unfiltered) frame the picture is clear
(Table~\ref{tab:coh-raw}): \texttt{normFalse} drops $35.7\%$ of pairs below
coherence $30$, against $0\%$ for \texttt{normTrue}.

\begin{table}[h]
  \centering
  \small
  \begin{tabular}{lccccc}
    \toprule
    scheme & $n$ & mean coh. & median & \% coh.${<}50$ & \% coh.${<}30$ \\
    \midrule
    \texttt{normFalse} & 28 & 44.2 & 40.9 & 60.7 & 35.7 \\
    \texttt{normTrue}  & 28 & 65.8 & 64.7 & 17.9 & 0.0  \\
    \texttt{per\_axis} & 28 & 41.9 & 39.4 & 75.0 & 21.4 \\
    \bottomrule
  \end{tabular}
  \caption{Steered coherence on the raw (unfiltered) frame,
  \texttt{power\_seeking}-excluded.}
  \label{tab:coh-raw}
\end{table}

Restricting to high positive-cosine pairs ($c_{ij}>0.25$) localises the confound
(Table~\ref{tab:coh-hicos}): \texttt{normFalse} breaks on $72.7\%$ of exactly
the pairs that carry the RQ1 cosine signal, while \texttt{normTrue} never falls
below coherence $30$.

\begin{table}[h]
  \centering
  \small
  \begin{tabular}{lccccc}
    \toprule
    scheme & $n$ & mean coh. & median & \% coh.${<}50$ & \% coh.${<}30$ \\
    \midrule
    \texttt{normFalse} & 11 & 29.2 & 25.7 & 81.8 & 72.7 \\
    \texttt{normTrue}  & 11 & 57.3 & 56.0 & 27.3 & 0.0  \\
    \texttt{per\_axis} & 11 & 43.6 & 41.7 & 81.8 & 9.1  \\
    \bottomrule
  \end{tabular}
  \caption{Steered coherence on high positive-cosine pairs ($c_{ij}>0.25$),
  raw frame, \texttt{power\_seeking}-excluded.}
  \label{tab:coh-hicos}
\end{table}

A row-level coherence gate (Filter~3, applied next) cannot repair this. On the
gated frame the \texttt{normFalse} mean does not look alarming
(Table~\ref{tab:coh-filtered}): the gate has simply discarded the broken
generations, so an average computed on the filtered file alone would understate
the problem.

\begin{table}[h]
  \centering
  \small
  \begin{tabular}{lcccc}
    \toprule
    scheme & $n$ & mean coh. & min & max \\
    \midrule
    \texttt{normFalse} & 18 & 56.1 & 31.8 & 86.2 \\
    \texttt{normTrue}  & 28 & 65.8 & 40.5 & 92.1 \\
    \texttt{per\_axis} & 22 & 47.1 & 31.3 & 83.3 \\
    \bottomrule
  \end{tabular}
  \caption{Steered coherence on the filtered ($\mathrm{coh}\geq 30$) frame,
  \texttt{power\_seeking}-excluded. The gate masks the damage.}
  \label{tab:coh-filtered}
\end{table}

The gate would remove $10$ of $28$ \texttt{normFalse} rows---the heaviest
censoring of any scheme---and those removals fall on the high-cosine pairs, so
the bias is in \emph{which} rows survive, correlated with the very predictor of
the RQ1 analysis. A row-level gate cannot fix a scheme-level confound;
\texttt{normFalse} is therefore dropped as a scheme \emph{before} gating and
treated as a diagnostic failure mode rather than a confirmatory result. The
normalized schemes (\texttt{normTrue}, \texttt{per\_axis}) hold the injection
magnitude independent of $c_{ij}$ and are the valid bases for the analysis.

\paragraph{Filter 3 (coherence gate): keep pair--scheme rows with mean coherence
at least $30$.}
With \texttt{normFalse} removed, the coherence gate is applied at the
pair--scheme row level to the two surviving schemes. For each joint condition we
average coherence over its $100$ generations and retain the row if the mean
coherence is at least $30$; this is the row-level step that produces the
\texttt{coh30} frame. The threshold was chosen after auditing generations below
$30$, where broken or repetitive text caused behaviour over-attribution by the
judge. A stricter threshold of $50$ discards too much readable signal, and
threshold-sensitivity checks preserve the qualitative directional result. Its
effect by scheme is summarised in Table~\ref{tab:coh-gate-counts}.

\begin{table}[h]
  \centering
  \small
  \begin{tabular}{lccc}
    \toprule
    scheme & raw dyads & retained & removed \\
    \midrule
    \texttt{normTrue} (normalized sum)         & 28 & 28 & 0 \\
    \texttt{per\_axis} (projection-controlled) & 28 & 22 & 6 \\
    \bottomrule
  \end{tabular}
  \caption{Effect of the coherence gate by scheme, on the two valid schemes.
  \texttt{normTrue} loses no dyads; \texttt{per\_axis} loses six.}
  \label{tab:coh-gate-counts}
\end{table}

Under normalized-sum injection (\texttt{normTrue}) no dyads are removed. Under
projection-controlled injection (\texttt{per\_axis}) six of $28$ dyads are
removed (Table~\ref{tab:coh-gate-censored}).

\begin{table}[h]
  \centering
  \small
  \begin{tabular}{lcc}
    \toprule
    pair & cosine & mean joint coherence \\
    \midrule
    formality + humorous       & $-0.52$ & $6.28$  \\
    apathetic + hallucinating  & $-0.16$ & $29.93$ \\
    confidence + humorous      & $-0.08$ & $26.45$ \\
    hallucinating + impolite   & $-0.05$ & $23.34$ \\
    hallucinating + humorous   & $+0.13$ & $23.18$ \\
    evil + humorous            & $+0.37$ & $28.32$ \\
    \bottomrule
  \end{tabular}
  \caption{The six \texttt{per\_axis} dyads censored by the coherence gate. Five
  of the six have cosine at most $+0.13$, consistent with the total-norm blow-up
  of projection-controlled injection for opposed directions.}
  \label{tab:coh-gate-censored}
\end{table}

\paragraph{Summary.}
\textbf{Filter 1 (behaviour) --- drop \texttt{power\_seeking}.} A trait-validity
filter, not a coherence filter: the vector is non-expressing (negative
single-trait $\Delta$), so every pair containing it has a structurally dead axis
and cannot inform the joint-vs-single regime question. The
\texttt{power\_seeking} pairs have high coherence (${\sim}86$), confirming the
exclusion is unrelated to coherence. This filter accounts for the $36$-vs-$28$
raw pair-count difference.

\textbf{Filter 2 (scheme) --- drop \texttt{normFalse}.} Its injection norm scales
with pair cosine ($\lVert\delta\rVert=\alpha\sqrt{2+2c_{ij}}$), so its coherence
collapse is concentrated on high-cosine pairs---exactly where the RQ1 cosine
signal lives. Because the degradation is correlated with the predictor, it
biases the regime comparison, and a row-level gate cannot fix a scheme-level
confound, so the scheme is dropped before gating. On the apples-to-apples
(\texttt{power\_seeking}-excluded) $28$-pair design, \texttt{normFalse} raw
coherence is the worst of the three schemes. \texttt{normTrue} and
\texttt{per\_axis} hold injection magnitude independent of cosine and are the
valid bases for the analysis.

\textbf{Filter 3 (coherence gate) --- keep rows with mean coherence at least
$30$.} The row-level gate that produces the \texttt{coh30} frame, applied last to
the two surviving schemes. Below $30$ the generations are broken or repetitive
and the judge over-attributes behaviour; a $50$ floor would discard too much
readable signal. It removes no \texttt{normTrue} dyads and six \texttt{per\_axis}
dyads (mostly antipodal, where projection-controlled injection's total norm blows
up).
